\section{An explicit cubic witness and a quantitative refinement}
\label{sec:plethysm-quantitative}

\paragraph{Status and dependence.}
The cubic fixed-locus and residue statements below are inherited from
ALPOGE\_GCT, P-F19, P-F23--P-F31, and P-F56.  The finite four-variable
calibration is ALPOGE\_GCT, P-F35.  We do not claim those results as new.
The revised dyadic estimates in PNEW-002--PNEW-005 are new quantitative
consequences of the same residue product; their proofs are given here.
All indices use the constituentwise conventions of GCT VI,
Definitions 1.2.2 and 1.2.4, printed pp.\ 13--14.

For $n\ge4$, put $g=\gcd(n,3)$, $d_0=n/g$, $M_n=2^{n-1}$,
$m_n=\binom{n+2}{3}$, and
\[
 f_n(s)=\dim\operatorname{Hom}_{GL_n}
  ((\det)^{3s/g},\operatorname{Sym}^{d_0s}(\operatorname{Sym}^3\mathbb C^n)).
\]
Set $w_i=(-2)^i$ for $0\le i<n-1$, and $w_{n-1}=0$.
Let $\mathcal S_n$ be the multiset consisting of the $n-1$ weights
$w_0,\ldots,w_{n-2}$ and all distinct-index triple sums
$w_i+w_j+w_k$, except the triples $\{0,1,k\}$ ($2\le k<n$)
and $\{0,2,3\}$.  Thus $|\mathcal S_n|=\binom n3$.
For a primitive $M_n$th root $\zeta$, the inherited exact residue is
\[
 C_n(\zeta)=\lim_{t\to\zeta}(1-t^{M_n})H_n(t)
       =\prod_{b\in\mathcal S_n}(1-\zeta^{gb})^{-1}.
 \tag{PNEW-R}
\]
Every factor is nonzero.  Define
\[
 P_n(a)=\prod_{b\in\mathcal S_n}|1-e^{2\pi iab/M_n}|^{-1},
 \qquad b_n=3^{-1}\pmod {M_n},
 \qquad
 Q_n=\sum_{\substack{a\bmod M_n\text{ odd}\\a\ne\pm b_n}}
            \frac{P_n(a)}{P_n(b_n)}.
\]
This is exactly $E_n/|Z_n|$ in P-F56: multiplication by the odd
number $g$ permutes primitive-root indices.  For $n\ge5$, the
increment $I_n(a)$ is the reciprocal product on
\[
 \mathcal T_n=\{1\}\sqcup
 \{1+w_j+w_k:2\le j<k<n,\ (j,k)\ne(2,3)\}
 \sqcup\{2+w_k:3\le k<n\}\sqcup\{6\}.
\]
Its cardinality is $\binom{n-1}{2}$, and cancellation of the shifted
smaller multiset gives
\[
 P_n(a)=I_n(a)P_{n-1}(-a\bmod M_{n-1}).
 \tag{PNEW-rec}
\]

\begin{proposition}[PNEW-001: a smallest-rank explicit witness]
For the cubic family one has
\[
                  2W_4(19)=\frac18>0=f_4(19).
\]
Consequently $p(f_4)>19$ and $s(f_4)\ge19$.  In fact the inherited
exact calibration gives $p(f_4)=s(f_4)=23$.  Rank four is the least
rank at which the primitive-$2^{n-1}$ wave gives an obstruction
$2W_n(x)>f_n(x)$ with $x\ge0$.
\end{proposition}
\begin{proof}
Salmon's cubic-surface invariant-ring theorem gives
\[
 H_4(t)=\frac{1+t^{25}}{(1-t^2)(1-t^4)(1-t^6)(1-t^8)(1-t^{10})}.
\]
See Huybrechts, \emph{The geometry of cubic hypersurfaces}, Chapter 4,
Theorem 4.1, p.\ 196, and the proof of Corollary 4.2; the complete
constituent calculation is ALPOGE\_GCT P-F35.  In
$\mathbb Q[z]/(z^4+1)$ the residue at a primitive eighth root is
\[
 C_4(z)=\frac{1+z^{25}}{(1-z^2)(1-z^4)(1-z^6)(1-z^{10})}
       =\frac{1+z+z^2+z^3}{8}.
\]
The primitive-dyadic trace formula, P-F56z, says that
$2W_4(x)=1/8$ for $x\equiv0,1,2,3\pmod8$, and $-1/8$ otherwise.
The numerator of $H_4$ has no odd term below degree 25, so
$f_4(19)=0$.  The strict-index conclusion is P-F56e'.

For $n=3$, the ternary-cubic invariant ring has generators of
ordinary degrees 4 and 6, so $H_3(t)=1/((1-t^4)(1-t^6))$.
An exact primary locator is McCrory--Shifrin--Varley (1994),
Section 3.2, printed p.~16.
At the primitive fourth roots $C_3=1/2$, hence
$2W_3(x)=\tfrac12\cos(\pi x/2)$.  Its positive values occur only
when $4\mid x$, where $f_3(x)\ge1$.  At $n=2$ one has
$H_2(t)=1/(1-t^2)$.  An elementary character check is as follows.
Put $d(j)=\#\{(a,b)\in\mathbb Z_{\ge0}^2:2a+3b=j\}$, with
$d(j)=0$ for $j<0$.  The difference of the weight-zero and
weight-two multiplicities in $\operatorname{Sym}^{2s}
(\operatorname{Sym}^3\mathbb C^2)$ is, from its Gaussian
polynomial character,
\[
 f_2(s)=d(3s)-d(s-1)-d(s-2)-d(s-3).
\]
The first term is $\lfloor s/2\rfloor+1$.  The sum of the last
three is $\lceil s/2\rceil$, since its generating function is
$(t+t^2+t^3)/((1-t^2)(1-t^3))=t/((1-t)(1-t^2))$.
Thus $f_2(s)$ equals one for even $s$ and zero for odd $s$, and
$2W_2(x)=(-1)^x\le f_2(x)$.  Rank one has no nontrivial primitive
dyadic subgroup.  These classical smaller-rank descriptions show
minimality; no asymptotic claim is inferred from the finite witness.
\end{proof}

\begin{lemma}[PNEW-002: a sharper near-increment estimate]
Let $n\ge5$, let $3\le j<M_n$ be odd, put
$a_j\equiv j3^{-1}\pmod{M_n}$, $t=\log_2j$, and
$R=\lceil\log_2(16j)\rceil$.  With
$B=(n-1)(n-4)/2$ one has
\[
 \left|\log_2\frac{I_n(a_j)}{I_n(b_n)}
        +Bt-\mathbf1_{3\mid j}\binom n2\right|
              \le \frac{16}{3}nR^2.
 \tag{PNEW-near}
\]
\end{lemma}
\begin{proof}
Write $l=n-1$.  The $B$ weights divisible by three are the active
weights $1+w_u+w_v$, $2\le u<v<l$, excluding $(2,3)$; the weights
$2+w_k$ for $3\le k<l$; and 6.  The remaining $l$ weights are
$1,2$, and $1+w_k$ for $2\le k<l$.  Give a weight its largest
active dyadic exponent $k$ and set $q=l-k$.  The exceptional
weights have one occurrence of each $q=1,\ldots,l$.
At each $q$ there are at most $n$ divisible weights: for a largest
index $k$ there are at most $k-2$ bulk pairs, one additional weight
$2+w_k$, and the separately assigned weight 6 at $k=2$.

Here are the elementary bounds needed for the comparison.  A sum
of at most $r\le3$ distinct signed dyadic terms has phase
$\kappa/3+(j/3)\sum_{i=1}^r\epsilon_i2^{-q_i}$, with
$q_1<\cdots<q_r$ and $\epsilon_i\in\{\pm1\}$.  Stop the initial
block before the first gap exceeding $R$.  Its numerator over
$3\cdot2^{q_u}$ is odd, so its distance from an integer is at least
$1/(3\cdot2^{q_u})$.  The discarded tail has magnitude at most
$1/(48\cdot2^{q_u})$, because $2^R\ge16j$.  Therefore the full
distance is at least $1/(6\cdot2^{q_1+(r-1)R})$.
Using $|1-e^{2\pi iy}|\ge4\|y\|$, a factor with $q_1\le R$ is
at least $2^{-3R-1}$; for two terms it is at least $2^{-2R-1}$.

Call $q>R$ low.  On a low divisible weight the cube-root shift
vanishes, and the target phase is $j$ times the reference phase.
Its magnitude is at most $(2j/3)2^{-q}<1/24$.
The inequalities $1-y^2/6\le\sin(y)/y\le1$ give error at most
$16j^2 4^{-q}$ from the main logarithm $t$.  Summing over at most
$n$ weights at each $q$ gives total error at most $n/48$.
For a high divisible weight the reference factor is at least
$2^{-q}/3$.  The preceding target bound and the upper bound 2
show that the absolute logarithm of the factor ratio is at most
$3R+2$.  Restoring its main term costs at most $t\le R$.
There are at most $nR$ high factors.  Thus the divisible-weight
error is bounded by
\[
                        4nR^2+2nR+n/48.
\]

Every exceptional reference factor lies in $[1/3,2]$.  If
$3\nmid j$, a low target factor has a nontrivial cube-root shift
perturbed by less than $1/24$, so its logarithmic ratio has
absolute value at most 3.  There are at most
$H=\min(R,l)$ high exceptional factors, each with logarithmic
ratio of absolute value at most $3R$.  Their total error is at
most $3n+3nR$.

If $3\mid j$, a low exceptional factor has logarithmic
reciprocal ratio $q-t$ with error at most 4.  Restoring the
missing high $q$ terms costs at most $H(H+1)/2\le nR$;
the high logarithmic ratios cost at most $3nR$; discarding
the low $-t$ terms costs at most $nR$.  Hence the error from
the exceptional main term $\sum_{q=1}^lq=\binom n2$ is at
most $5nR+4n$.

Since $j\ge3$, $R\ge6$.  In the second, larger case the total
error divided by $nR^2$ is at most
\[
                4+\frac7R+\frac{193}{48R^2}
       \le4+\frac76+\frac{193}{1728}<\frac{16}{3}.
\]
The first case is smaller.  This proves the lemma with all
constants uniform in $n$ and $j$.
\end{proof}

\begin{lemma}[PNEW-003: a sharper far-increment estimate]
For odd $a\bmod M_n$, put $j\equiv3a\pmod{M_n}$,
$1\le j<M_n$, and $J=\min(j,M_n-j)$.  Then
\[
 \log_2\frac{I_n(a)}{I_n(b_n)}
       \le \frac{25}{24}n^2+3n-\frac{(\log_2J-8)_+^3}{24}.
 \tag{PNEW-far}
\]
\end{lemma}
\begin{proof}
Set $l=n-1$, $\theta=a/2^l$, and
$y_{uv}=\theta(1+w_u+w_v)$ for $2\le u<v<l$.
Distinct subsets of the active dyadic weights have distinct sums
modulo $2^l$: their difference has a unique smallest dyadic
valuation.  This remains true after multiplication by odd $a$.
For $0\le h<l$, put $r=l-h$ and count the pairs with
$\|y_{uv}\|<2^{-h}$.  There are $e=(r-2)_+$ early indices.
Early--early pairs number at most $\binom e2$.  For each early
index, good late partners lie in one residue class modulo $2^r$;
the open interval $(-2^r,2^r)$ contains at most two representatives
of that class.  Good late--late pairs similarly number at most
two.  Thus the layer count is at most
$\binom e2+2e+2$, which is $\binom r2+1$ for $r\ge2$ and 2
for $r=1$.  Its sum over all layers is $\binom n3+n$.

If $J>3\cdot2^r$, pair the early indices consecutively into
$\lfloor e/2\rfloor$ disjoint pairs.  For two such pairs,
\[
 y_{uv}+y_{u+1,v}+y_{u,v+1}=3\theta\pmod1.
\]
At least one of the four associated cross factors is outside the
current layer.  Distinct blocks use disjoint factors, and the
omitted factor $(2,3)$ is internal to an index-pair.  The layer
count therefore loses at least $\binom{\lfloor e/2\rfloor}{2}$.
Summing for $r\le(\lfloor\log_2J\rfloor-2)_+$ gives loss at
least $(\log_2J-8)_+^3/24$, by summing
$\binom{\lfloor(r-2)/2\rfloor}{2}\ge(r-5)^2/8$ for $r\ge6$.

The remaining increment factors are exactly the two families
\[
       \{c\}\sqcup\{c+w_k:2\le k<l\},\qquad c=1,2.
\]
In each layer there are at most $(r-2)_+$ early terms and at
most two late terms, including the singleton $c$ in the late
class.  Each family's total layer count is at most
$\binom n2+1$.

For every nonzero dyadic phase $y$,
\[
 \log_2|1-e^{2\pi iy}|^{-1}
   \le -\log_2\|y\|-2
   \le \sum_{h=0}^{l-1}\mathbf1_{\{\|y\|<2^{-h}\}}-2.
\]
Add the omitted bulk factor only to the nonnegative layer counts,
and subtract 2 for each of the actual $\binom{n-1}{2}$ increment
factors.  The resulting upper bound is
\[
 \log_2 I_n(a)\le\binom n3+3n
                    -\frac{(\log_2J-8)_+^3}{24}.
\]
For the selected increment, all $B=(n-1)(n-4)/2$ weights
divisible by three can be estimated at a small phase.  Write
$\alpha=13/12$.  The classical inequality $\pi<22/7$ and
the exact integer inequality $(44/21)^{12}<2^{13}$ give
$2\pi/3<2^\alpha$.  For an active bulk weight with largest
index $k$ and other index $j$, its selected factor is at most
\[
  2^\alpha\frac{2^k}{M_n}
                       (1+2^{j-k}+2^{-k}).
\]
For $2+w_k$, $3\le k<n-1$, replace the parenthesis by
$1+2^{1-k}$; for the remaining weight 6, assign $k=2$ and
use the parenthesis $3/2$.  The inequality
$\log_2(1+u)\le2u$ for $u\ge0$ follows from
$\ln(1+u)\le u$ and $\ln2>1/2$.
For the bulk corrections,
$\sum_{j=2}^{k-1}2^{j-k}<1$ and
$\sum_{k=3}^\infty(k-2)2^{-k}=1/2$ give a total at most
$2(n-5)+1$; omitting the pair $(2,3)$ only improves it.
The $2+w_k$ corrections sum to at most 1, and the correction
for 6 is less than 1.  Their combined sum is therefore at
most $2n$ for every $n\ge5$.

The sum of $n-1-k$ over all divisible weights is exactly
\[
 \binom{n-2}{3}-(n-4)+\binom{n-3}{2}+(n-3)
                  =\binom{n-1}{3}-n+4.
\]
Each of the $n-1$ exceptional factors is at most 2.  Hence
\[
 \log_2 I_n(b_n)\ge
       \binom{n-1}{3}-\frac{13}{12}B-4n+5.
\]
Subtracting this from the preceding layer bound leaves
$25n^2/24+67n/24-11/6$, which is at most
$25n^2/24+3n$.  This proves the lemma.
\end{proof}

\begin{lemma}[PNEW-004: a cubic bound for the starting mass]
For every $n\ge5$,
\[
                  Q_n\le2^{\,2n^3/3+10n/3-4}.
 \tag{PNEW-mass}
\]
\end{lemma}
\begin{proof}
Use layers as in the preceding proof, now with active indices
$0,\ldots,n-2$.  For distinct active triple sums, the layer
count is at most $\binom r3+2\binom r2+2r+2$, by fixing the
early subset and using at most two late representatives in each
residue class.  For distinct active pair sums it is at most
$\binom r2+2r+2$.  For singleton weights it is at most $r$:
a nonzero late multiple of $2^r$ cannot lie in the open interval
$(-2^r,2^r)$.  Add omitted weights only to the nonnegative layer
counts.  Subtracting 2 for the actual $\binom n3$ factors gives
\[
 \log_2P_n(a)\le\binom n4+\binom n3+5\binom n2+4(n-1).
\]
Here is a direct selected-amplitude estimate from the exact product,
avoiding any further determinant formula.  Among all active distinct
triples $0\le i<j<k\le n-2$, each weight is divisible by three and
has absolute value less than $2^{k+1}$.  Its selected factor is at
most $8\cdot2^k/M_n$, using $\pi<4$.  Summing $n-1-k$
over all such triples gives $\binom n4$.  To obtain the actual
bulk in $\mathcal S_n$, delete $\{0,1,k\}$ for $2\le k\le n-2$
and $\{0,2,3\}$.  The sum of their contributions $n-1-k-3$
is $(n-3)(n-8)/2+(n-7)$.  Bound each remaining nonbulk factor
by 2.  There are $\binom n3$ actual factors altogether.  Thus
\[
 \begin{split}
 \log_2P_n(b_n)&\ge
 \binom n4-2\binom{n-1}{3}-\binom n3
                                  -\frac{n^2}{2}+\frac{7n}{2}-3\\
 &=\frac{n^4}{24}-\frac{3n^3}{4}
                        +\frac{59n^2}{24}-\frac{3n}{4}-1.
 \end{split}
 \tag{PNEW-amplitude}
\]
There are fewer than $2^{n-1}$ summands in $Q_n$.  Combining
these bounds yields
\[
                  \log_2Q_n\le\frac23n^3+\frac{10}{3}n-4.
\]
If $Q_n=0$, the asserted inequality is immediate.
\end{proof}

\begin{theorem}[PNEW-005: an explicit uniform threshold]
For every integer $n\ge79000$,
\[
                  p(f_n)>2^{n/8},\qquad s(f_n)\ge2^{n/8}.
\]
Thus the inherited plethysm refutations of PH2 and SH(a) have an
explicit all-rank threshold of $79000$ by elementary estimates.
This is a sufficient threshold, not a claim of optimality.
\end{theorem}
\begin{proof}
Put $N=39500$ and $c=1168/12500$.  We first show that for every $n\ge N$ and
every nonselected odd $a\bmod M_n$,
\[
                   I_n(a)/I_n(b_n)\le2^{-2n^2/7}.
\]
By conjugation use the least positive absolute residue
$J$ of $3a$ and set $t=\log_2J$.  If $J\ne1$, then
$t>19/12$, since $3^{12}>2^{19}$.  For $t\le cn$, PNEW-002 and $R\le t+5$ give
the upper bound
\[
 F_n(t)=-\frac{(t-1)n^2}{2}
          +\frac52nt+\frac{16}{3}n(t+5)^2.
\]
This convex quadratic takes its maximum on
$[19/12,cn]$ at an endpoint.  The first endpoint gives
\[
 F_n(19/12)=-\frac{7n^2}{24}+\frac{50783}{216}n.
\]
After division by $n^2$, both endpoint upper bounds decrease
with $n$.  At $N$ their margins below $-2/7$ are exactly
\[
 -\frac27-\frac{F_N(19/12)}{N^2}
       =\frac{19}{59724000}>0,\qquad
 -\frac27-\frac{F_N(cN)}{N^2}
       =\frac{8303507}{86406250}>0.
\]
For $t\ge cn$, PNEW-003 and the exact rational inequality
\[
 \frac{(cN-8)^3}{24N^2}-\frac{25}{24}-\frac3N-\frac27
       =\frac{1676026865671}{255978515625000}>0
\]
give the same bound.  The left side is increasing with $N$
replaced by any real $n\ge N$: the derivative of
$(cn-8)^3/n^2$ is $(cn-8)^2(cn+16)/n^3>0$,
and $-3/n$ is increasing.

Grouping the two lifts of every nonselected parent in
PNEW-rec gives the exact recurrence
\[
 Q_n\le2^{1-2n^2/7}Q_{n-1}+2\delta_n,\qquad
 \delta_n\le2^{-n^3/12}\quad(n\ge48),
\]
where the last bound is the selected-sibling product P-F56ae.
In particular every contraction factor is at most $1/2$ for
$n\ge N$.  PNEW-004 and iteration therefore give
\[
 Q_n\le
 2^{(2/3)N^3+(10/3)N+(n-N+1)-2(n^3-(N-1)^3)/21}
                  +4\cdot2^{-N^3/12}.
\]
We used
$\sum_{k=N}^n k^2\ge\int_{N-1}^{n}x^2\,dx$.
At $n=2N$ the first exponent is exactly
$-445603261$, hence less than $-3$; it decreases for
all larger $n$.  The second summand is less than $1/8$.
Thus $Q_n<1/4$ for every $n\ge2N=79000$.

For completeness, the conversion to an index bound uses the
following amplitude and phase facts.  If $r=3g\in\{3,9\}$ and
$a_*r\equiv1\pmod{M_n}$, then
\[
 |Z_n|\ge2^{n^4/24-2n^3}.
\]
Indeed $|Z_n|=P_n(b_n)/M_n$, and PNEW-amplitude is stronger.
Among any $r$ consecutive integer arguments, the selected
conjugate pair has a value at least $|Z_n|/2$: its phases differ
from $r$ equally spaced phases by at most $2\pi/M_n$, and
$\cos(\pi/r)\ge1/2$.  For $n\ge6$ the error is at most
$1/4$.  Subtracting the remaining mass $Q_n|Z_n|\le|Z_n|/4$
gives an $x$ in
$[\lceil2^{n/8}\rceil,\lceil2^{n/8}\rceil+8]$ with
\[
                    2W_n(x)\ge2^{n^4/24-2n^3-1}.
\]
On the other hand
$f_n(x)\le\binom{m_n+d_0x-1}{m_n-1}$.
For $n\ge1024$, use $m_n\le n^3/5$,
$\log_2n\le n/16$, $2^{n/8}\ge n^2$, and
$x\le2^{n/8+1}$ to obtain
\[
 \log_2\max(1,f_n(x))
 \le\frac{n^3}{5}(n/8+\log_2n+2)
 \le\frac{3n^4}{80}+\frac{2n^3}{5}
                  <\frac{n^4}{24}-2n^3-1.
\]
The last difference is $n^4/240-12n^3/5-1>0$.
Hence $2W_n(x)>f_n(x)$.  The exact paired-constituent identity
$q_a(X)-q_{a+L/2}(X)=2W_n(a)$ implies that any positivity shift
at most $x$, or saturation shift below $x$, is impossible.
This proves the theorem for the displayed sufficient threshold.
\end{proof}
